\documentclass{dlproblemset}

\title{Funções de Custo e Funções de Ativação}
\subtitle{Lista de Exercícios}

\begin{document}

\begin{enumerate}[1)]

    \item Dentre as funções de ativação abaixo, qual é a opção mais adequada para modelar um problema de classificação \textbf{multiclasse} no qual você quer interpretar a saída da rede como as probabilidades da observação pertencer a uma dada classe?
    \begin{enumerate}[(a)]
        \item ReLU: $f(x_{i}) = max(0,x_{i})$
        \item Sigmoid: $f(x_{i}) = \frac{1}{1 + e^{-x_{i}}}$
        \item Tanh: $f(x_{i}) = \frac{2}{1 + e^{-2x_{i}}}-1$
        \item Softplus: $f(x_{i}) = \log_{e}(1 + e^{x_{i}})$
        \item Softmax $f(x_{i}) = \frac{e^{x_{i}}}{\sum_{z} e^{x_{z}}}$
        \item Linear $f(x_{i}) = x_{i}$
    \end{enumerate}
    %e

    \item Dentre as funções de ativação abaixo, qual é a opção mais adequada para modelar um problema de classificação \textbf{multilabel} no qual você quer interpretar a saída da rede como as probabilidades da observação pertencer a uma dada classe?
    \begin{enumerate}[(a)]
        \item ReLU: $f(x_{i}) = max(0,x_{i})$
        \item Sigmoid: $f(x_{i}) = \frac{1}{1 + e^{-x_{i}}}$
        \item Tanh: $f(x_{i}) = \frac{2}{1 + e^{-2x_{i}}}-1$
        \item Softplus: $f(x_{i}) = \log_{e}(1 + e^{x_{i}})$
        \item Softmax $f(x_{i}) = \frac{e^{x_{i}}}{\sum_{z} e^{x_{z}}}$
        \item Linear $f(x_{i}) = x_{i}$
    \end{enumerate}
    %b

    \item ``Para um problema de classificação \textbf{multilabel} onde a saída pode ter 10 classes diferentes (simultaneamente ou não) a função de ativação na camada de saída do MLP deve ser \_\_\_\_\_\_\_  e o treinamento deve buscar otimizar a função de perda \_\_\_\_\_\_\_''. Quais opções abaixo melhor preenchem as lacunas da afirmação acima?

    \begin{enumerate}[(a)]
        \item ReLU: $f(x_{i}) = max(0,x_{i})$
        \item Sigmoid: $f(x_{i}) = \frac{1}{1 + e^{-x_{i}}}$
        \item Tanh: $f(x_{i}) = \frac{2}{1 + e^{-2x_{i}}}-1$
        \item Softplus: $f(x_{i}) = \log_{e}(1 + e^{x_{i}})$
        \item Softmax $f(x_{i}) = \frac{e^{x_{i}}}{\sum_{z} e^{x_{z}}}$
        \item Linear $f(x_{i}) = x_{i}$
	    \item \textit{Binary Cross Entropy Loss}: $L_{\text{bce}} = -(ylog(\hat{y}) + (1 - y)log(1-\hat{y}))$.
	    \item \textit{Categorical Cross Entropy Loss}: $L_{\text{cce}} = - \sum_{i=1}^{c}y_{i}\log\hat{y}_{i}$
	    \item \textit{Sum of Squared Residuals}: $L_{\text{ssr}} = (y - \hat{y})^{2}$.
	    \item \textit{Absolute Error Loss}: $L_{\text{abs}} = \left |y - \hat{y}\right|$.
	    \item \textit{Zero-One Loss}: $L_{01} = I(y \neq \hat{y})$.
    \end{enumerate}
    %b, g

    \item Durante o treinamento de uma rede neural você percebe que os gradientes diminuem rapidamente até que se aproximam de zero, fazendo que os pesos da rede permaneçam praticamente sem mudança. Você rapidamente identifica que se trata de um caso de \textit{vanishing gradient}, marque dentre as opções abaixo possíveis causas para esse problema.

	\begin{enumerate}[(a)]
	    \item As camadas ocultas usam função de ativação \textit{Tanh}.
	    \item As camadas ocultas usam função de ativação \textit{ReLU}.
	    \item As camadas ocultas usam função de ativação \textit{Sigmoid}.
	    \item A rede usa \textit{batch normalization}.
	\end{enumerate}
    %a, c

    \item A função de ativação \textit{Rectified Linear Unit (ReLU)} é uma das funções de ativação mais utilizadas em redes neurais. A \textit{ReLU} é definida como $f(x) = \max(0,x)$. Sobre essa função de ativação, é correto afirmar que:
    \begin{enumerate}[(a)]
	    \item Minimiza o problema de \textit{vanishing gradient} em comparação com a função de ativação sigmoid.
	    \item É continua e diferenciável em todo seu domínio.
	    \item Não é contínua e não é diferenciável em todo seu domínio.
	    \item É contínua e não é diferenciável em todo seu domínio.
    \end{enumerate}
    %a, d

    \item Dentre as opções abaixo, quais adicionam não linearidades na rede neural.

    \begin{enumerate}[(a)]
        \item Adicionar camadas de convolução.
        \item Diminuir a taxa de aprendizado.
        \item Embaralhar o conjunto de treinamento.
        \item Usar função de ativação ReLU.
        \item Adicionar regularização.
    \end{enumerate}
    %d

    \item O preço de um determinado veículo $i$ pode ser modelado satisfatoriamente pela seguinte função $f(\mathbf{X}^{(i)}) = 10000 X^{(i)}_{1} + 50000 X^{(i)}_{2} + 3500$. Responda as questões abaixo supondo que você está utilizando um \textit{MultiLayer Perceptron} para aproximar o preço do veículo e justifique sua resposta.
    \begin{enumerate}[(a)]
	    \item Quantos neurônios/unidades são necessários no mínimo para modelar o preço dos veículos?
	    \item Qual função de ativação você deveria utilizar?
    \end{enumerate}

    \item Cite e explique duas características desejáveis de uma função de ativação em redes neurais.

    \item Por que é usual trabalhar com funções de ativação não-lineares em redes neurais?

    \item Calcule as derivadas das seguintes funções de ativação:
    \begin{enumerate}[a)]
        \item ReLU $\varphi(x)=max(0,x)$
        \item Sigmoid $\varphi(x)=\frac{1}{(1+e^{-x})}$
        \item \emoji{skull-and-crossbones} Softmax $\varphi(\mathbf{x})=S_i=\frac{e^{x_i}}{\Sigma e^{x_j}}$, (encontrar a matriz Jacobiana)
    \end{enumerate}

    \item Qual a utilidade das funções de ativação nas camadas ocultas de um MLP?

    \item Qual a utilidade das funções de ativação na camada de saída de um MLP?

    \item Calcule as derivadas das seguintes funções de custo:
    \begin{enumerate}[a)]
        \item  $L_{l2}=\sum\left(y^{(i)}-\hat{y}^{(i)}\right)^2$
        \item $L_{BCE}=-yln(\hat{y})-(1-y)ln(1-\hat{y})$
    \end{enumerate}

    \item Desenvolva a função de custo do Erro Médio Quadrático a partir da distribuição Gaussiana.
    \begin{equation*}
        Pr(y|\mu,\sigma^2)=\frac{1}{\sqrt{2\pi\sigma^2}}exp\left(-\frac{(y-\mu)^2}{2\sigma^2}\right)
    \end{equation*}

    \item Desenvolva a função de custo Entropia Binária Cruzada a partir da distribuição de Bernoulli.
    \begin{equation*}
        Pr(y|\lambda)= (1-\lambda)^{(1-y)}\lambda^y
    \end{equation*}

    \item \emoji{skull-and-crossbones} Suponha que você quer criar uma rede neural para estimar a direção $y$ (em radianos) do vento a partir de medições de pressão $\mathbf{x}$.
    Uma distribuição que trabalha no domínio circular é a distribuição de von Mises:
    \begin{equation*}
        Pr(y|\mu,k) = \frac{exp[k*cos(y-\mu)]}{2\pi * Bessel_0[k]}
    \end{equation*}
    \noindent onde $\mu$ é a direção média e $k$ é o inverso da variância. O termo $Bessel_0[k]$ é uma função modificada de Bessel com grau 0.
    Crie uma função de custo para aprender o parâmetro $\mu$ dessa distribuição.

    \item \emoji{skull-and-crossbones} Suponha que você quer criar uma rede neural para estimar o número de pedestres $y \in \{0,1,2,...\}$ que passam em uma rua da cidade em algum determinado momento. Você possui um vetor de atributos $\mathbf{x}$ que contém informações sobre a hora do dia, sobre o clima e sobre o dia ser feriado ou não.
    Uma distribuição de probabilidade adequada para trabalhar com esse problema é a distribuição de Poisson:
    \begin{equation*}
        Pr(y=k) = \frac{\lambda^ke^{-\lambda}}{k!}
    \end{equation*}
    \noindent onde $\lambda$ é o único parâmetro, que representa a média da distribuição.
    Crie uma função de custo assumindo que temos acesso a $I$ instâncias de treinamento compostas por pares $\{x^{i},y^{i}\}$.

    \item Sobre as asserções abaixo a respeito de funções de ativação:

    \textbf{I.} A função sigmoide $\sigma(z) = 1/(1 + e^{-z})$ satura para entradas com módulo grande, fazendo com que a sua derivada tenda a zero. Em redes profundas, isso contribui para o problema do desaparecimento do gradiente nas camadas iniciais.

    \textbf{II.} A função ReLU é estritamente positiva para qualquer entrada, garantindo que nenhuma unidade da rede possa ``morrer'' (deixar permanentemente de produzir saída diferente de zero) durante o treinamento.

    \textbf{III.} A derivada da função tanh tem magnitude máxima maior do que a da sigmoide ($\tanh'(0) = 1$ contra $\sigma'(0) = 0.25$), de modo que, em regiões não saturadas, o gradiente propagado é tipicamente maior, o que costuma permitir uma convergência mais rápida do treinamento em comparação com a sigmoide.

    \begin{enumerate}[(a)]
        \item Apenas I está correta.
        \item Apenas II está correta.
        \item Apenas III está correta.
        \item Apenas I e II estão corretas.
        \item Apenas I e III estão corretas.
        \item Apenas II e III estão corretas.
        \item Todas estão corretas.
        \item Nenhuma está correta.
    \end{enumerate}
    %e

    \item Você está treinando um MLP profundo com 10 camadas ocultas para um problema de classificação binária. Optou por sigmoide em todas as camadas ocultas, sigmoide na camada de saída, função de custo entropia cruzada binária, otimizador SGD com taxa de aprendizado padrão e inicialização $\mathcal{N}(0, 1)$ para todos os pesos. Após muitas épocas, o erro de treino permanece alto e os gradientes em relação aos pesos das primeiras camadas são extremamente próximos de zero. Qual a explicação mais consistente \textbf{e} a melhor mitigação?
    \begin{enumerate}[(a)]
        \item A taxa de aprendizado está baixa demais; aumentá-la em duas ordens de grandeza compensa os gradientes pequenos e resolve o problema.
        \item O modelo está sofrendo \textit{overfitting}; reduzir o número de camadas para uma única camada oculta resolve o problema.
        \item Os gradientes nas camadas iniciais ficam próximos de zero devido à saturação da sigmoide nas camadas profundas; trocar a sigmoide das ocultas por ReLU com inicialização de He resolve.
        \item A função de custo entropia cruzada binária é inadequada para classificação; substituí-la por L2 loss resolve o problema.
        \item Os pesos devem ser inicializados em zero para evitar que comecem em regiões saturadas da sigmoide.
        \item Trocar a sigmoide da camada de saída por SoftMax resolve o desaparecimento do gradiente nas camadas iniciais.
        \item Adicionar Dropout com probabilidade alta nas camadas iniciais aumenta o gradiente nessas camadas.
        \item O problema decorre exclusivamente do otimizador SGD; substituí-lo por Adam, sem outras alterações, é suficiente para que o erro caia.
    \end{enumerate}
    %c

    \item Por que, na prática, a função de ativação da camada de saída de uma rede neural é escolhida \emph{em conjunto} com a função de custo (por exemplo, sigmoide com entropia cruzada binária; SoftMax com entropia cruzada multinomial)?
    \begin{enumerate}[(a)]
        \item Porque é a única forma de evitar \textit{overfitting} na camada de saída.
        \item Porque sem essa combinação a rede é incapaz, em princípio, de produzir saídas no intervalo correto (por exemplo, $[0, 1]$ para sigmoide).
        \item Porque o algoritmo de \textit{backpropagation} só funciona quando função de custo e ativação da saída são compostas algebricamente.
        \item Porque o teorema da aproximação universal exige especificamente esse pareamento para garantir convergência da descida de gradiente.
        \item Porque essa combinação é o único caso em que a função de custo é convexa em relação aos pesos da rede como um todo.
        \item Porque essa combinação simplifica o gradiente em relação aos \textit{logits}, resultando em $\partial \mathcal{L} / \partial z = \hat{y} - y$.
        \item Porque essa combinação reduz o número de parâmetros treináveis na última camada.
        \item Porque, sem ela, a derivada da função de custo em relação aos pesos da última camada é sempre zero.
    \end{enumerate}
    %f

    \item Avalie as cinco afirmações a seguir sobre funções de custo, verossimilhança e ativações:

    \textbf{I.} A entropia cruzada binária corresponde à máxima verossimilhança sob a hipótese $y \mid \mathbf{x} \sim \mathrm{Bernoulli}(\hat{y})$, enquanto o erro quadrático médio corresponde à máxima verossimilhança sob hipótese Gaussiana com variância fixa.

    \textbf{II.} A derivada da sigmoide atinge seu valor máximo em $z = 0$, com $\sigma'(0) = \tfrac{1}{4}$, enquanto a da tangente hiperbólica atinge $\tanh'(0) = 1$.

    \textbf{III.} Sendo $S_i = \mathrm{softmax}(\mathbf{z})_i$, a matriz Jacobiana da \textit{softmax} em relação aos seus \textit{logits} é diagonal, isto é, $\partial S_i / \partial z_j = 0$ para todo $i \neq j$.

    \textbf{IV.} O erro quadrático médio não pode, em hipótese alguma, ser usado como função de custo em classificação binária: a rede não treina.

    \textbf{V.} A ReLU satisfaz $\mathrm{ReLU}'(z) = 1$ para $z > 0$ e $\mathrm{ReLU}'(z) = 0$ para $z < 0$; portanto, ao contrário da sigmoide, não satura para entradas positivas grandes, o que ajuda a mitigar o desaparecimento do gradiente.

    \begin{enumerate}[(a)]
        \item Apenas I, II e V estão corretas.
        \item Apenas I e II estão corretas.
        \item Apenas I, II, IV e V estão corretas.
        \item Apenas II, III e V estão corretas.
        \item Apenas I, II e III estão corretas.
        \item Apenas I, IV e V estão corretas.
        \item Todas as afirmações estão corretas.
        \item Apenas III e IV estão corretas.
    \end{enumerate}
    %a

    \item Marque \textbf{Verdadeiro} ou \textbf{Falso} nas asserções abaixo. Em caso \textbf{Falso}, proponha uma modificação \emph{não-trivial} que torna a asserção verdadeira (a simples negação \textbf{não} é uma modificação válida).
    \begin{enumerate}[a)]
        \item ( )\,V \quad ( )\,F \quad Um MLP com $1$ unidade de entrada, $K$ camadas ocultas com $D$ unidades cada (todas com viés) e $1$ unidade de saída possui exatamente $3D + 1 + (K - 1)\,D^2$ parâmetros treináveis.
        \item ( )\,V \quad ( )\,F \quad A descida de gradiente em lote (\textit{batch gradient descent}) calcula o gradiente sobre uma única instância (ou um \textit{mini-batch} pequeno) a cada iteração; já a SGD (\textit{Stochastic Gradient Descent}) calcula o gradiente da função de custo sobre todo o conjunto de treinamento a cada iteração.
        \item ( )\,V \quad ( )\,F \quad O problema do desaparecimento do gradiente é particularmente severo com a função de ativação sigmoide em redes profundas, pois a derivada $\sigma'(z) = \sigma(z)(1-\sigma(z))$ tem valor máximo $0.25$ e tende a zero quando a sigmoide satura (i.e., para $|z|$ grande), fazendo com que o gradiente decaia multiplicativamente ao longo das camadas.
        \item ( )\,V \quad ( )\,F \quad A função de custo de entropia cruzada binária pode ser obtida a partir da estimação por máxima verossimilhança dos rótulos $y_i \in \{0, 1\}$ assumindo que cada rótulo segue uma distribuição \emph{Gaussiana} de média $\hat{y}_i$ e variância fixa.
        \item ( )\,V \quad ( )\,F \quad Inicializar todos os pesos de uma camada oculta com zero é adequado, pois o algoritmo de \textit{backpropagation} naturalmente quebra a simetria entre as unidades nas primeiras iterações da descida de gradiente.
    \end{enumerate}

    \item A distribuição paramétrica \textit{Krék Normal}, com parâmetro de posição $\mu$ e escala $\sigma$ fixa e conhecida (não predita pela rede), tem densidade
    \[
    \Pr(y \mid \mu, \sigma) \;\propto\;
    \exp\!\left(-\frac{(y - \mu)^2}{2\sigma^2}\right) \cdot
    \bigl|\cos(y - \mu)\bigr|^{\,1/(2\sigma^2)} .
    \]
    Você quer treinar uma rede que prediz $\mu^{(i)} = f(\mathbf{x}^{(i)}; \boldsymbol{\theta})$ para $N$ observações independentes, seguindo a receita de máxima verossimilhança vista em aula. Descartando tudo que não altera o ponto de mínimo em $\boldsymbol{\theta}$, qual é a função de custo a ser minimizada?

    \begin{enumerate}[(a)]
        \item $\mathcal{L} = \displaystyle\sum_{i=1}^{N} \left[ \bigl(y^{(i)} - \mu^{(i)}\bigr)^2 + \ln\bigl|\cos(y^{(i)} - \mu^{(i)})\bigr| \right]$
        \item $\mathcal{L} = \displaystyle\sum_{i=1}^{N} \left[ \bigl(y^{(i)} - \mu^{(i)}\bigr)^2 - \ln\bigl|\cos(y^{(i)} - \mu^{(i)})\bigr| \right]$
        \item $\mathcal{L} = \displaystyle\sum_{i=1}^{N} \left[ \frac{\bigl(y^{(i)} - \mu^{(i)}\bigr)^2}{2\sigma^2} - \ln\bigl|\cos(y^{(i)} - \mu^{(i)})\bigr| \right]$
        \item $\mathcal{L} = \displaystyle\sum_{i=1}^{N} \left[ -\bigl(y^{(i)} - \mu^{(i)}\bigr)^2 - \ln\bigl|\cos(y^{(i)} - \mu^{(i)})\bigr| \right]$
        \item $\mathcal{L} = \displaystyle\sum_{i=1}^{N} \left[ \bigl(y^{(i)} - \mu^{(i)}\bigr)^2 \cdot \ln\bigl|\cos(y^{(i)} - \mu^{(i)})\bigr| \right]$
        \item $\mathcal{L} = \displaystyle\sum_{i=1}^{N} \left[ \bigl(y^{(i)} - \mu^{(i)}\bigr)^2 - \bigl|\cos(y^{(i)} - \mu^{(i)})\bigr| \right]$
        \item $\mathcal{L} = \displaystyle\sum_{i=1}^{N} \left[ \bigl(y^{(i)} - \mu^{(i)}\bigr)^2 - \ln\bigl|\cos(\mu^{(i)})\bigr| \right]$
        \item $\mathcal{L} = -\displaystyle\sum_{i=1}^{N} \left[ \bigl(y^{(i)} - \mu^{(i)}\bigr)^2 - \ln\bigl|\cos(y^{(i)} - \mu^{(i)})\bigr| \right]$
    \end{enumerate}
    %b

    \item Associe cada escolha de função de ativação (coluna da esquerda) ao comportamento observado durante o treinamento (coluna da direita).

    \begin{center}
    \begin{tabular}{p{0.34\linewidth} p{0.58\linewidth}}
    \textbf{1.} Sigmoide em todas as camadas ocultas de uma rede profunda & \textbf{(i)} as ativações ficam centradas em zero, mas o gradiente ainda decai nas regiões em que as unidades operam saturadas. \\
    \textbf{2.} ReLU com vieses deslocados para valores muito negativos & \textbf{(ii)} parte das unidades passa a produzir saída zero para todas as instâncias e deixa de receber gradiente. \\
    \textbf{3.} Tangente hiperbólica em todas as camadas ocultas & \textbf{(iii)} o gradiente que chega às camadas iniciais é multiplicado por um fator de no máximo $0{,}25$ a cada camada. \\
    \textbf{4.} Identidade (linear) em todas as camadas ocultas & \textbf{(iv)} a rede inteira equivale a uma única transformação afim, qualquer que seja o número de camadas. \\
    \end{tabular}
    \end{center}

    \begin{enumerate}[(a)]
        \item 1-(i); 2-(ii); 3-(iii); 4-(iv)
        \item 1-(iii); 2-(i); 3-(ii); 4-(iv)
        \item 1-(ii); 2-(iii); 3-(i); 4-(iv)
        \item 1-(iii); 2-(ii); 3-(iv); 4-(i)
        \item 1-(i); 2-(iii); 3-(ii); 4-(iv)
        \item 1-(iv); 2-(ii); 3-(i); 4-(iii)
        \item 1-(ii); 2-(i); 3-(iii); 4-(iv)
        \item 1-(iii); 2-(ii); 3-(i); 4-(iv)
    \end{enumerate}
    %h

    \item Um colega propõe treinar um classificador binário linear com saída sigmoide usando a função de custo quadrática $J = \tfrac{1}{2}(y - \hat{y})^2$ no lugar da entropia cruzada binária. Sobre a superfície de custo resultante, qual afirmação é correta?
    \begin{enumerate}[(a)]
        \item O custo quadrático composto com a sigmoide permanece convexo nos parâmetros, apenas com o mínimo deslocado em relação ao da entropia cruzada.
        \item É a entropia cruzada que perde a convexidade quando composta com a sigmoide, enquanto o custo quadrático preserva essa propriedade.
        \item As duas funções perdem a convexidade quando compostas com a sigmoide, de modo que a escolha entre elas é indiferente para a otimização.
        \item A convexidade não depende da função de custo escolhida, mas apenas do número de camadas ocultas presentes na arquitetura da rede.
        \item O custo quadrático composto com a sigmoide deixa de ser convexo nos parâmetros, e a otimização pode parar em mínimos locais em vez do global.
        \item O custo quadrático deixa de ser diferenciável nos pontos em que a sigmoide satura, o que impede o uso de descida de gradiente no treinamento.
        \item O custo quadrático passa a ter derivada constante em todo o domínio, fazendo a descida de gradiente divergir para qualquer taxa de aprendizado.
        \item A perda de convexidade é irrelevante na prática, pois a descida de gradiente sempre encontra o mínimo global de qualquer função diferenciável.
    \end{enumerate}
    %e

    \item Você quer prever o número de chamadas que uma central telefônica recebe por hora, $y \in \{0, 1, 2, \ldots\}$, modelando a saída com a distribuição de Poisson
    \[
    \Pr(y \mid \lambda) = \frac{\lambda^{y}\,e^{-\lambda}}{y!}, \qquad \lambda > 0 ,
    \]
    em que a rede prediz $\lambda^{(i)} = f(\mathbf{x}^{(i)}; \boldsymbol{\theta})$ para cada uma das $N$ observações independentes. Seguindo a receita de máxima verossimilhança vista em aula e descartando tudo que não altera o ponto de mínimo em $\boldsymbol{\theta}$, qual combinação de função de custo e função de ativação na saída é adequada? As ativações consideradas nas alternativas são
    \[
    \sigma(z) = \frac{1}{1 + e^{-z}}, \qquad
    \mathrm{softplus}(z) = \ln\bigl(1 + e^{z}\bigr), \qquad
    \text{identidade: } f(z) = z, \qquad
    \mathrm{ReLU}(z) = \max(0, z).
    \]

    \begin{enumerate}[(a)]
        \item $\mathcal{L} = \sum_{i}\bigl[\lambda^{(i)} - y^{(i)}\ln\lambda^{(i)}\bigr]$, com sigmoide na camada de saída.
        \item $\mathcal{L} = \sum_{i}\bigl[\lambda^{(i)} - y^{(i)}\ln\lambda^{(i)}\bigr]$, com \textit{softplus} na camada de saída.
        \item $\mathcal{L} = \sum_{i}\bigl[\lambda^{(i)} - y^{(i)}\ln\lambda^{(i)}\bigr]$, com identidade (linear) na camada de saída.
        \item $\mathcal{L} = \sum_{i}\bigl[\lambda^{(i)} - y^{(i)}\ln\lambda^{(i)}\bigr]$, com ReLU na camada de saída.
        \item $\mathcal{L} = \sum_{i}\bigl[y^{(i)}\ln\lambda^{(i)} - \lambda^{(i)}\bigr]$, com \textit{softplus} na camada de saída.
        \item $\mathcal{L} = \sum_{i}\bigl[y^{(i)}\ln\lambda^{(i)} - \lambda^{(i)}\bigr]$, com sigmoide na camada de saída.
        \item $\mathcal{L} = \sum_{i}\bigl[y^{(i)}\ln\lambda^{(i)} - \lambda^{(i)}\bigr]$, com ReLU na camada de saída.
        \item $\mathcal{L} = \sum_{i}\bigl[y^{(i)}\ln\lambda^{(i)} - \lambda^{(i)}\bigr]$, com identidade (linear) na camada de saída.
    \end{enumerate}
    %b

\end{enumerate}

\newpage

\subsection*{Gabarito}
\color{blue}
\begin{enumerate}[1)]

    \item e. A softmax normaliza os \textit{logits} para que as saídas sejam positivas e somem 1, produzindo uma distribuição de probabilidade sobre classes mutuamente exclusivas, que é o que o problema multiclasse exige. As demais ativações produzem valores por unidade sem a restrição de somarem 1, então não formam uma distribuição sobre as classes.

    \item b. No problema \textit{multilabel} cada classe é uma decisão binária independente, então cada saída precisa de uma probabilidade própria em $[0,1]$, o que a sigmoid fornece. A softmax (e) forçaria as classes a competir entre si, pois as probabilidades somariam 1, o que impede múltiplas classes simultâneas.

    \item b, g. Sigmoid em cada unidade de saída fornece probabilidades independentes por classe, e a \textit{Binary Cross Entropy} avalia cada uma dessas decisões binárias separadamente. A combinação softmax + \textit{Categorical Cross Entropy} seria a escolha para o caso multiclasse com classes mutuamente exclusivas, que não é o cenário do enunciado.

    \item a, c. Tanh e sigmoid saturam para entradas de magnitude grande e têm derivadas menores que 1 (a da sigmoid vale no máximo $0.25$); no \textit{backpropagation} essas derivadas se multiplicam camada após camada e o gradiente encolhe exponencialmente. A ReLU (b) tem derivada 1 na região ativa, e o \textit{batch normalization} (d) é uma técnica que combate o problema em vez de causá-lo.

    \item a, d. Na região ativa ($x>0$) a derivada da ReLU é constante igual a 1, então ela não satura como a sigmoid e minimiza o \textit{vanishing gradient} (a). Ela é contínua em todo o domínio, mas não é diferenciável em $x=0$, onde as derivadas laterais diferem (d), o que exclui (b) e (c).

    \item d. A ReLU é uma função não linear aplicada às pré-ativações, e é isso que impede a rede de colapsar em um modelo linear. Convoluções (a) são operações lineares, e taxa de aprendizado (b), embaralhamento dos dados (c) e regularização (e) afetam o treinamento, não a classe de funções que a rede consegue representar.

    \item
    \begin{enumerate}[(a)]
        \item Uma única unidade é suficiente. A função alvo é linear nos atributos, e um neurônio computa exatamente uma combinação linear das entradas mais um viés: basta fazer $\theta_1 = 10000$, $\theta_2 = 50000$ e $\theta_0 = 3500$.
        \item A ativação linear (identidade). O problema é de regressão e a saída deve poder assumir qualquer valor real; uma ativação não linear na saída restringiria ou distorceria a imagem da função modelada.
    \end{enumerate}

    \item Exemplos de características desejáveis (citar e explicar duas): ser não linear, pois é o que permite que a composição de camadas represente funções que um modelo linear não representa; ser diferenciável (ou quase em todo ponto), pois o treinamento depende da propagação de gradientes; não saturar em grande parte do domínio, pois derivadas próximas de zero causam \textit{vanishing gradient}; ter baixo custo computacional, pois a ativação é avaliada em todas as unidades a cada \textit{forward} e \textit{backward}.

    \item Porque sem não linearidade entre as camadas a composição de transformações lineares colapsa em uma única transformação linear ($\mathbf{\Theta}^{(2)}(\mathbf{\Theta}^{(1)} \mathbf{X}) = (\mathbf{\Theta}^{(2)}\mathbf{\Theta}^{(1)})\mathbf{X}$). Nesse caso a rede, por mais profunda que seja, teria o mesmo poder de representação de um modelo linear. As ativações não lineares quebram essa equivalência e permitem aproximar funções complexas.

    \item
    \begin{enumerate}[a)]
        \item \[
\varphi'(x) =
     \begin{cases}
       \text{0,} &\quad x \le0,\\
       \text{1,} &\quad x > 0\\
     \end{cases}
\]
        \item $\varphi'(x)=\varphi(x)(1-\varphi(x))$
        \item \emoji{skull-and-crossbones} \[
        \frac{\partial S_i}{\partial x_j} =
     \begin{cases}
       S_i(1-S_i), &\quad i=j,\\
       -S_i S_j,   &\quad i\neq j.
     \end{cases}
\]
        Forma compacta: $\dfrac{\partial S_i}{\partial x_j} = S_i\,(\delta_{ij} - S_j)$, onde $\delta_{ij}$ é o delta de Kronecker.
    \end{enumerate}

    \item Adicionar não linearidades. Sem elas, a composição das camadas colapsaria em uma única transformação linear e a profundidade não acrescentaria poder de representação; com ativações não lineares nas camadas ocultas, cada camada pode compor funções progressivamente mais complexas.

    \item Ajustar a saída da rede neural à imagem da função modelada: a ativação da última camada mapeia as pré-ativações para o domínio que a tarefa exige, como $[0,1]$ com a sigmoid para probabilidade binária, uma distribuição que soma 1 com a softmax no caso multiclasse, ou os reais com a ativação linear em regressão. Essa escolha também casa a saída com a função de custo correspondente.

    \item
    \begin{enumerate}[a)]
    \item $2(\hat{y}-y)$
    \item $\frac{\hat{y}-y}{\hat{y}(1-\hat{y})}$ \quad ({\footnotesize esta é $\partial L_{BCE} / \partial \hat y$, não $\partial L_{BCE} / \partial \theta$; ao multiplicar pela derivada da sigmoid $\hat y(1-\hat y)$ e por $\mathbf{X}^\top$, simplifica para $\mathbf{X}^\top(\hat{\mathbf{y}} - \mathbf{y})$})
    \end{enumerate}

    \item Aplicando o logaritmo negativo à verossimilhança da Gaussiana com $\mu = \hat{y}^{(i)}$: $-\ln Pr(y^{(i)}|\hat{y}^{(i)},\sigma^2) = \frac{1}{2}\ln(2\pi\sigma^2) + \frac{(y^{(i)}-\hat{y}^{(i)})^2}{2\sigma^2}$. O primeiro termo e o fator $\frac{1}{2\sigma^2}$ não dependem de $\hat{y}$, então maximizar a verossimilhança equivale a minimizar $L_{l2}=\sum\left(y^{(i)}-\hat{y}^{(i)}\right)^2$.

    \item Aplicando o logaritmo negativo à verossimilhança da Bernoulli com $\lambda = \hat{y}$: $-\ln\left[(1-\hat{y})^{(1-y)}\,\hat{y}^{\,y}\right] = -y\,ln(\hat{y})-(1-y)\,ln(1-\hat{y}) = L_{BCE}$.

    \item \emoji{skull-and-crossbones} $-\sum_{i} cos(y^{(i)}-f(x^{(i)},\theta))$ \quad ({\footnotesize a NLL é $-\sum_i \ln \Pr(y^{(i)}\mid \mu, k)$. O denominador $2\pi \cdot \mathrm{Bessel}_0[k]$ e o $k$ não dependem de $\mu = f(x^{(i)},\theta)$, então caem como constantes; sobra apenas o termo $-k\sum_i \cos(y^{(i)} - f(x^{(i)},\theta))$, que com $k$ fixo é proporcional ao apresentado})

    \item \emoji{skull-and-crossbones} $\sum _{i=1}^{I}(f[x^{(i)},\theta] - y^{(i)}ln(f[x^{(i)},\theta]))$ \quad ({\footnotesize a NLL de cada observação é $-\ln Pr(y^{(i)}) = -y^{(i)}\ln\lambda + \lambda + \ln(y^{(i)}!)$; com $\lambda = f[x^{(i)},\theta]$, o termo $\ln(y^{(i)}!)$ não depende de $\theta$ e cai como constante, restando a soma apresentada})

    \item e. \textbf{I é correta:} a saturação da sigmoide leva $\sigma'(z) \to 0$ para $|z|$ grande; em redes profundas o produto de derivadas pequenas decai ao longo das camadas (gradiente desaparece nas iniciais). \textbf{II é incorreta:} a ReLU \textit{pode} morrer (\textit{dying ReLU}). Se uma unidade passa a receber consistentemente entradas negativas, sua saída fica em zero e $\mathrm{ReLU}'(z) = 0$, de modo que o gradiente em relação aos pesos dessa unidade é nulo e ela deixa de aprender. \textbf{III é correta:} $\tanh'(z) = 1 - \tanh^2(z)$ atinge máximo $1$ em $z = 0$, enquanto $\sigma'(z) = \sigma(z)(1 - \sigma(z))$ atinge máximo $0.25$ em $z = 0$. Em regiões não saturadas a tanh propaga gradientes até cerca de $4\times$ maiores do que a sigmoide, o que tipicamente acelera a convergência da descida de gradiente em comparação com a sigmoide. Logo, apenas I e III estão corretas. As alternativas (a), (b) e (c) excluem afirmações corretas; (d), (f) e (g) incluem a afirmação II, que é falsa; (h) contradiz I e III.

    \item c. A combinação ``sigmoide em todas as camadas + rede profunda + inicialização $\mathcal{N}(0, 1)$'' produz pré-ativações com módulo grande, levando a saturação da sigmoide; $\sigma'$ pequeno multiplicado camada a camada faz com que o gradiente desapareça nas camadas iniciais. A correção padrão é trocar a sigmoide das ocultas por ReLU/Leaky ReLU (cujas derivadas não saturam) e usar inicialização adequada (He), mantendo sigmoide na saída por se tratar de classificação binária. Em (a), aumentar a taxa de aprendizado em duas ordens de grandeza tipicamente leva a divergência, e não cura gradientes que já chegam praticamente nulos às primeiras camadas devido à saturação da sigmoide; em (b), reduzir a profundidade ``resolve'' o sintoma mas remove a vantagem representacional da rede, portanto não é a melhor mitigação; em (d), L2 loss combinada com sigmoide na saída \emph{piora} o gradiente (introduz fator $\sigma'$ adicional); em (e), pesos zerados eliminam toda quebra de simetria e a rede não consegue aprender; em (f), SoftMax é apropriada para multiclasse com rótulo único, o problema é binário e a saturação ocorre nas camadas \emph{ocultas}, não na saída; em (g), \textit{dropout} desliga unidades e não aumenta o gradiente; em (h), Adam ajusta passos por estimativas adaptativas, mas se o gradiente real é praticamente zero o sinal de aprendizado continua ausente: Adam ajuda, mas isolado não cura o problema da saturação.

    \item f. Com sigmoide + entropia cruzada binária (e analogamente SoftMax + entropia cruzada multinomial), o gradiente $\partial \mathcal{L}/\partial z$ em relação aos \textit{logits} simplifica-se a $\hat{y} - y$. Esse cancelamento elimina o termo $\sigma'(z)$, que tende a zero quando a sigmoide satura, e fornece um sinal de aprendizado bem-comportado mesmo quando a predição está muito errada, fator crítico para a estabilidade do treinamento. Em (a), \textit{overfitting} é tratado por regularização, não pela escolha do par ativação/custo da saída; (b) é parcialmente verdadeira para o intervalo, mas \emph{não} é a razão principal, pois outras combinações também produziriam saídas no intervalo, sem o ganho computacional; em (c), \textit{backpropagation} funciona com qualquer composição diferenciável; em (d), o teorema da aproximação universal não exige nenhum pareamento específico de saída/custo; em (e), a função de custo da rede como um todo \emph{não} é convexa nos pesos, mesmo com esse pareamento (a convexidade vale apenas para a regressão logística com uma única camada); em (g), a escolha não altera a contagem de parâmetros; em (h), sem o pareamento o gradiente \emph{não} é zero, é apenas pior comportado (com termos de saturação).

    \item a. \textbf{I (V):} BCE deriva da máxima verossimilhança Bernoulli; MSE da Gaussiana com variância fixa. \textbf{II (V):} $\sigma'(z)$ é máxima em $z = 0$, com $\sigma'(0) = \tfrac{1}{4}$, e $\tanh'(0) = 1$. \textbf{III (F):} a matriz Jacobiana da \textit{softmax} \emph{não} é diagonal: $\partial S_i/\partial z_j = -S_i S_j$ para $i \neq j$ (e $S_i(1-S_i)$ para $i = j$). \textbf{IV (F):} o MSE \emph{pode} ser usado em classificação binária (a rede treina), apenas é menos bem fundamentado estatisticamente e costuma convergir mais devagar do que a entropia cruzada; ``não treina'' é falso. \textbf{V (V):} $\mathrm{ReLU}'(z) = 1$ para $z > 0$, sem saturação no ramo positivo, mitigando o desaparecimento do gradiente. Verdadeiras: I, II e V. A alternativa (e) inclui III (falsa); (c) e (f) incluem IV (falsa); (b) omite V; (d) troca I por III; (g) e (h) incluem afirmações falsas.

    \item
    \begin{enumerate}[a)]
        \item \textbf{F.} A expressão $(K - 1)\,D^2$ conta apenas os pesos das $K - 1$ transições entre camadas ocultas, ignorando os $D$ \textit{biases} de cada camada oculta de destino (ou seja, das camadas $h_2, h_3, \ldots, h_K$), totalizando $(K - 1)\,D$ parâmetros omitidos. Decompondo (pesos + \textit{biases} da camada de destino): entrada para a primeira oculta tem $1 \cdot D + D = 2D$ parâmetros; cada uma das $K - 1$ transições entre camadas ocultas tem $D \cdot D + D = D(D+1)$ parâmetros (não apenas $D^2$); última oculta para a saída tem $D \cdot 1 + 1 = D + 1$ parâmetros. \textit{Modificação:} substituir $(K - 1)\,D^2$ por $(K - 1)\,D(D+1)$, totalizando $3D + 1 + (K - 1)\,D(D+1)$.
        \item \textbf{F.} As definições estão invertidas. \textit{Modificação:} a \textbf{descida de gradiente em lote} (\textit{batch GD}) calcula o gradiente sobre \emph{todo} o conjunto de treinamento a cada iteração; a \textbf{SGD} (\textit{Stochastic Gradient Descent}) calcula o gradiente sobre uma \emph{única} instância (ou um \textit{mini-batch} pequeno) a cada iteração.
        \item \textbf{V.} Descrição correta do desaparecimento do gradiente em redes profundas com sigmoide. Também vale mencionar o efeito multiplicativo $(\le 0.25)^L$ ao longo de $L$ camadas.
        \item \textbf{F.} A entropia cruzada binária decorre da máxima verossimilhança assumindo $y_i \mid \hat{y}_i \sim \mathrm{Bernoulli}(\hat{y}_i)$, \emph{não} Gaussiana. Para $y \in \{0,1\}$, a verossimilhança é $p(y \mid \hat{y}) = \hat{y}^{y}(1-\hat{y})^{1-y}$ e $-\log p$ resulta exatamente em $-[y\log \hat{y} + (1-y)\log(1-\hat{y})]$. A hipótese Gaussiana com variância fixa, ao contrário, conduz ao MSE (regressão), não à entropia cruzada. \textit{Modificação:} substituir ``Gaussiana de média $\hat{y}_i$ e variância fixa'' por ``Bernoulli de parâmetro $\hat{y}_i$''.
        \item \textbf{F.} Com pesos zero, todas as unidades de uma camada computam a mesma pré-ativação, recebem o mesmo gradiente e portanto sofrem a mesma atualização: a simetria \emph{nunca} é quebrada por \textit{backpropagation}, e as unidades permanecem idênticas indefinidamente. \textit{Modificação:} é necessário inicializar os pesos com pequenos valores \emph{aleatórios} (e.g., Xavier/He); só essa aleatoriedade quebra a simetria entre unidades de uma mesma camada.
    \end{enumerate}

    \item b. Seguindo a receita vista em aula, a função de custo é a \textit{negative log-likelihood}. Para observações independentes a verossimilhança é o produto das densidades, e o logaritmo transforma produto em soma. Note que o expoente $1/(2\sigma^2)$ incide sobre \emph{os dois} fatores da densidade:
    \[
    -\ln \Pr(y \mid \mu, \sigma)
     = \frac{(y - \mu)^2}{2\sigma^2} - \frac{1}{2\sigma^2}\ln\bigl|\cos(y - \mu)\bigr| + \ln Z(\sigma)
     = \frac{1}{2\sigma^2}\Bigl[ (y - \mu)^2 - \ln\bigl|\cos(y - \mu)\bigr| \Bigr] + \ln Z(\sigma),
    \]
    \[
    \mathcal{L} = \frac{1}{2\sigma^2}\sum_{i=1}^{N} \Bigl[ \bigl(y^{(i)} - \mu^{(i)}\bigr)^2 - \ln\bigl|\cos(y^{(i)} - \mu^{(i)})\bigr| \Bigr] + N\ln Z(\sigma).
    \]
    Duas coisas podem ser descartadas, por motivos diferentes. A constante de normalização $N \ln Z(\sigma)$ é um termo \textbf{aditivo} que não depende de $\boldsymbol{\theta}$: ela desloca o valor da loss sem mover o mínimo. Já $1/(2\sigma^2)$ é um fator \textbf{multiplicativo global e positivo}, que multiplica tudo que depende de $\boldsymbol{\theta}$; como $\arg\min_{\boldsymbol{\theta}} c\,f(\boldsymbol{\theta}) = \arg\min_{\boldsymbol{\theta}} f(\boldsymbol{\theta})$ para $c > 0$, ele também sai sem alterar a solução. Sobra
    \[
    \mathcal{L} = \sum_{i=1}^{N} \Bigl[ \bigl(y^{(i)} - \mu^{(i)}\bigr)^2 - \ln\bigl|\cos(y^{(i)} - \mu^{(i)})\bigr| \Bigr].
    \]
    Lendo o resultado: o primeiro termo é o erro quadrático usual, herdado do fator exponencial. O segundo depende apenas do resíduo $r^{(i)} = y^{(i)} - \mu^{(i)}$: como $-\ln|\cos(r)|$ vale $0$ em $r = 0$ e cresce sem limite quando $|r| \to \pi/2$, ele penaliza resíduos grandes de forma bem mais agressiva que a quadrática nessa faixa. A alternativa (a) não nega o segundo termo: minimizar $+\ln|\cos(r)|$ empurra o resíduo \emph{para longe} de zero, na direção oposta à indicada pela verossimilhança. (c) é o erro mais instrutivo: mantém $1/(2\sigma^2)$ apenas no termo quadrático, como se o expoente da densidade incidisse só no fator exponencial; aí o $\sigma$ deixa de ser um fator global e passa a ser o \emph{peso relativo} entre os dois termos, e o ponto de mínimo muda com $\sigma$ (um fator constante só pode ser descartado quando multiplica a expressão inteira). (d) nega apenas o termo quadrático, o que faria a rede maximizar o erro de predição. (e) mantém o produto dos dois fatores; o logaritmo da verossimilhança converte produto em soma, não em produto. (f) esquece o logaritmo sobre $|\cos(r)|$, mantendo esse fator na escala original enquanto o outro termo já está no domínio logarítmico. (g) aplica o cosseno a $\mu$ em vez do resíduo $y - \mu$, trocando um termo que depende do erro por um que depende só da predição. (h) inverte o sinal da soma inteira, o que corresponde a \emph{maximizar} a NLL em vez de minimizá-la.

    \item h. \textbf{1-(iii):} a derivada da sigmoide, $\sigma'(z) = \sigma(z)(1 - \sigma(z))$, tem máximo $0{,}25$ em $z = 0$. Empilhando camadas, a regra da cadeia multiplica esses fatores e o gradiente decai exponencialmente com a profundidade, mesmo sem saturação. \textbf{2-(ii):} vieses muito negativos deixam a pré-ativação abaixo de zero para todas as instâncias. Como $\mathrm{ReLU}'(z) = 0$ para $z < 0$, a unidade nunca recebe gradiente e não consegue voltar, o que caracteriza o \textit{dying ReLU}. \textbf{3-(i):} a tanh é centrada em zero, o que a diferencia da sigmoide e ajuda na propagação, mas ela também satura nos dois extremos, com $\tanh'(z) \to 0$ para $|z|$ grande. \textbf{4-(iv):} a composição de transformações afins é uma transformação afim. Sem não-linearidade,
    \[
    \mathbf{\Theta}^{(2)}\bigl(\mathbf{\Theta}^{(1)}\mathbf{x} + \mathbf{b}^{(1)}\bigr) + \mathbf{b}^{(2)}
    = \bigl(\mathbf{\Theta}^{(2)}\mathbf{\Theta}^{(1)}\bigr)\mathbf{x} + \bigl(\mathbf{\Theta}^{(2)}\mathbf{b}^{(1)} + \mathbf{b}^{(2)}\bigr),
    \]
    que é uma única camada, independentemente da profundidade. Em (a), sigmoide e tanh foram trocadas: o fator $0{,}25$ é o máximo da derivada da sigmoide, não da tanh; (b) dá à ReLU a saturação centrada em zero da tanh e à tanh o \textit{dying ReLU}, que exige derivada exatamente nula em um intervalo; (c) dá o \textit{dying ReLU} à sigmoide, cuja derivada nunca é exatamente zero, e o fator $0{,}25$ à ReLU; (d) acerta as duas primeiras associações, mas troca tanh e identidade, e a tanh não colapsa a rede em uma transformação afim; (e) tira o fator $0{,}25$ da sigmoide e coloca o \textit{dying ReLU} na tanh, que satura sem nunca zerar a derivada; (f) troca sigmoide e identidade, e a sigmoide não torna a rede equivalente a uma única camada; (g) embaralha as três primeiras descrições, deixando o \textit{dying ReLU} na sigmoide e o fator $0{,}25$ na tanh.

    \item e. Em um classificador linear o \textit{logit} é $z = \mathbf{x}^{\top}\boldsymbol{\theta}$, afim nos parâmetros, de modo que a convexidade da função de custo em $z$ se transfere para os parâmetros. A entropia cruzada binária é convexa em $z$; já o erro quadrático composto com a sigmoide, $\bigl(\sigma(z) - y\bigr)^2$, não é. Como a sigmoide satura nas duas pontas, essa composição produz regiões praticamente planas e curvatura que troca de sinal, e a superfície passa a admitir mínimos locais e platôs em que a descida de gradiente pode estacionar longe do mínimo global. Foi por isso que, ao adaptar a regressão linear para classificação, trocamos o custo quadrático pela entropia cruzada. A alternativa (a) afirma que a convexidade é preservada, o que é justamente o que se perde na composição com a sigmoide; (b) inverte os papéis, pois a entropia cruzada é a que permanece convexa em $z$; (c) supõe que ambas perdem a convexidade, o que tornaria a escolha indiferente, mas a entropia cruzada não perde; (d) atribui a convexidade à profundidade da rede, mas aqui o modelo é linear e a diferença vem exclusivamente da função de custo; (f) confunde perda de convexidade com perda de diferenciabilidade, e a sigmoide é diferenciável em todo o domínio, inclusive nas regiões saturadas; em (g), a derivada não é constante, ela é próxima de zero nas regiões saturadas, o que trava o treinamento em vez de fazê-lo divergir; em (h), a descida de gradiente não garante o mínimo global de funções não convexas, que é exatamente o problema levantado.

    \item b. A NLL das $N$ observações é
    \[
    -\sum_{i}\ln\Pr\bigl(y^{(i)} \mid \lambda^{(i)}\bigr) = \sum_{i}\Bigl[\lambda^{(i)} - y^{(i)}\ln\lambda^{(i)} + \ln y^{(i)}!\Bigr],
    \]
    e $\ln y^{(i)}!$ não depende de $\boldsymbol{\theta}$, logo $\mathcal{L} = \sum_i\bigl[\lambda^{(i)} - y^{(i)}\ln\lambda^{(i)}\bigr]$. A saída da rede precisa ser estritamente positiva e sem limite superior, porque $\lambda$ é a média de uma contagem; a \textit{softplus}, $\ln(1 + e^{z})$, atende as duas exigências. Em (a), a sigmoide limita $\lambda$ a $(0, 1)$ e impede prever médias acima de uma chamada por hora; em (c), a identidade permite $\lambda \leq 0$ e deixa $\ln\lambda$ indefinido; em (d), a ReLU produz $\lambda = 0$ em toda a região negativa, $\ln 0$ diverge e o gradiente se anula; (e) escreve a log-verossimilhança sem o sinal negativo, e minimizá-la equivale a \emph{minimizar} a verossimilhança; (f), (g) e (h) combinam o sinal trocado de (e) com as ativações inadequadas de (a), (d) e (c), respectivamente.
\end{enumerate}

\end{document}
